\begin{solapartat}
\punts{0,15} $^{222}_{86}\mathrm{Rn} \rightarrow\ ^{218}_{84}\mathrm{Po} + ^{4}_{2}\alpha$

\punts{0,15} $^{218}_{84}\mathrm{Po} \rightarrow\ ^{214}_{82}\mathrm{Pb} + ^{4}_{2}\alpha$

\punts{0,15} $^{214}_{82}\mathrm{Pb} \rightarrow\ ^{214}_{83}\mathrm{Bi} + ^{0}_{-1}e + \bar{\nu}_e$

\punts{0,15} $^{214}_{83}\mathrm{Bi} \rightarrow\ ^{214}_{84}\mathrm{Po} + ^{0}_{-1}e + \bar{\nu}_e$

\punts{0,15} $^{214}_{84}\mathrm{Po} \rightarrow\ ^{210}_{82}\mathrm{Pb} + ^{4}_{2}\alpha$

\punts{0,15} $^{210}_{82}\mathrm{Pb} \rightarrow\ ^{210}_{83}\mathrm{Bi} + ^{0}_{-1}e + \bar{\nu}_e$

\punts{0,15} $^{210}_{83}\mathrm{Bi} \rightarrow\ ^{210}_{84}\mathrm{Po} + ^{0}_{-1}e + \bar{\nu}_e$

\punts{0,2} $^{210}_{84}\mathrm{Po} \rightarrow\ ^{206}_{82}\mathrm{Pb} + ^{4}_{2}\alpha$

\destaca{Alternativament} es pot escriure $^{4}_{2}\mathrm{He}$ en lloc de $^{4}_{2}\alpha$, o bé $^{0}_{-1}\beta^{-}$ en lloc de $^{0}_{-1}e^{-}$.
\end{solapartat}

\begin{solapartat}
\punts{0,75} El període de semidesintegració correspon al temps necessari per desintegrar la meitat dels nuclis inicials, $8\cdot 10^{23}/2 = 4\cdot 10^{23}$.

\figura[0.55]{sol_fig1.png}

Del gràfic, obtenim $T_{1/2} = 3,5$ dies, és a dir, corresponen a la desintegració del $^{210}_{83}\mathrm{Bi}$.

\punts{0,25} $N(t) = N_0\,e^{-\lambda t}$
\[ \frac{1}{2} = e^{-\lambda T_{1/2}} \Rightarrow \lambda = -\frac{\ln 0,5}{T_{1/2}} = 0,200\ \mathrm{dies^{-1}} \]

\punts{0,25} $N(t) = 8\times 10^{23} - 7,95\times 10^{23} = 5,0\times 10^{21}$
\[ \frac{N(t)}{N_0} = \frac{5\times 10^{21}}{8\times 10^{23}} = 6,25\times 10^{-3} = e^{-\lambda t} \Rightarrow t = -\frac{\ln(6,25\times 10^{-3})}{\lambda} = 25,4\ \text{dies} \]

\destaca{Alternativament}

\punts{0,25} $N(t) = 8\times 10^{23} - 7,95\times 10^{23} = 5,0\times 10^{21}$
\[ \frac{N(t)}{N_0} = \frac{5\times 10^{21}}{8\times 10^{23}} = 6,25\times 10^{-3} \]

\punts{0,25}
\[ \left.\begin{aligned} 6,25\times 10^{-3} &= e^{-\lambda t}\\ 0,5 &= e^{-\lambda T_{1/2}} \end{aligned}\right\} \Rightarrow \frac{\ln 6,25\times 10^{-3}}{\ln 0,5} = \frac{t}{T_{1/2}} \]
\[ t = T_{1/2}\,\frac{\ln 6,25\times 10^{-3}}{\ln 0,5} = 25,4\ \text{dies} \]
\end{solapartat}
